\chapter{Donanim Guvenligi}

\section*{Giris}
Donanim guvenligi, sistemin guven kokunu olusturur. Isletim sistemi yuklenmeden once calisan firmware ve cipsel guvenlik mekanizmalari, tum ust katmanlarin guvenilirligini belirler.

\section{Cip (Chipset) ve Anakart Seviyesinde Guvenlik}
\subsection{Guvenilir Platform Modulu (TPM) ve Kriptografik Islemciler}
TPM, anahtar koruma, olculmus onyukleme ve uzaktan butunluk dogrulama icin temel bilesendir.
\begin{itemize}
    \item Disk sifreleme anahtarlarinin guvenli saklanmasi
    \item Cihaz saglik kaniti (attestation)
    \item Kimlik tabanli erisim politikalariyla entegrasyon
\end{itemize}

\section{Firmware ve BIOS/UEFI Guvenligi}
\subsection{Guvenli Onyukleme (Secure Boot) Surecleri}
UEFI Secure Boot, yalnizca imzali bileşenlerin yuklenmesini saglar. Uretim ortami sistemlerinde firmware guncelleme zinciri dogrulanabilir olmalidir.
\begin{enumerate}
    \item Guvenli anahtar hiyerarsisi kur
    \item Imzali bootloader ve kernel dagit
    \item Firmware degisikliklerini denetim kaydina al
\end{enumerate}

\section{Donanim Tedarik Zinciri Riskleri}
\subsection{Sahte Bilesenler ve Donanimsal Truva Atlari (Hardware Trojans)}
Tedarik zinciri guvenligi, sadece satin alma degil; kabul testi, seri no dogrulama, firmware hash kontrolu ve tedarikci denetimini kapsar.

\Not{Kritik sistemler icin tedarik zinciri risk puanlamasi yapilmali; tek tedarikci bagimliligini azaltan ikincil kaynak plani korunmalidir.}
